\documentclass[12pt]{article}
\usepackage[T1]{fontenc}
\usepackage[utf8]{inputenc}
\usepackage{lmodern}
\usepackage{textcomp}
\usepackage{enumitem}
\usepackage{geometry}
\geometry{margin=1in}
\begin{document}
\begin{center}
Answer Key - Geography - K-12 Level Assessment 24 \
\small (Answers are ordered exactly as in the question paper.)
\end{center}

\section*{Multiple Choice}
\begin{enumerate}[label=\arabic*.]
\item \textbf{Answer:} A
\\ \textit{Options:} A) First; B) Second; C) Third; D) Fourth
\\ \textit{Note:} The first zone in Burgess's concentric zone model, known as the "Zone of Transition," is characterized by low-income slums and ethnic ghettos due to its proximity to industrial areas and urban decay.
\item \textbf{Answer:} D
\\ \textit{Options:} A) Southeast Asia; B) Sub-Saharan Africa; C) North America; D) Latin America
\\ \textit{Note:} Spanish colonial architecture, uneven economic development, and rural-to-urban migration flows are characteristic features of Latin America due to its historical colonization by Spain and the resulting socio-economic dynamics.
\item \textbf{Answer:} D
\\ \textit{Options:} A) Transnational corporations; B) EU; C) World Trade Organization; D) Good transportation network
\\ \textit{Note:} A good transportation network facilitates trade, movement, and communication, thereby strengthening the state's functionality rather than presenting a challenge.
\item \textbf{Answer:} D
\\ \textit{Options:} A) European; B) Muslim; C) Sub-Saharan African; D) Latin American
\\ \textit{Note:} The Latin American city model features a distinct residential spine that radiates outward from the central business district along major boulevards, reflecting the urban layout and socio-economic patterns typical of many Latin American cities.
\item \textbf{Answer:} D
\\ \textit{Options:} A) Vietnam; B) United Kingdom; C) Germany; D) Bolivia
\\ \textit{Note:} Bolivia does not have a well-known example of a relict boundary, as its territorial delineations are primarily based on natural features and political agreements rather than on historical remnants of past boundaries that no longer serve their original purpose.
\end{enumerate}

\section*{True / False}
\begin{enumerate}[label=\arabic*.]
\setcounter{enumi}{5}
\item \textbf{Answer:} False
\\ \textit{Options:} A) True; B) False
\\ \textit{Note:} Paris, France, is the capital city and is located within the central core area of the country, making the statement false.
\item \textbf{Answer:} False
\\ \textit{Options:} A) True; B) False
\\ \textit{Note:} The statement is false because the European Union does not have a single common language; instead, it recognizes 24 official languages, allowing for multilingual communication among its member states.
\item \textbf{Answer:} False
\\ \textit{Options:} A) True; B) False
\\ \textit{Note:} The statement is false because structural assimilation refers to the integration of immigrants into the social and economic institutions of the host society, while cultural assimilation involves the merging of cultural identities, which does not necessarily occur in the case of immigrants learning English.
\item \textbf{Answer:} False
\\ \textit{Options:} A) True; B) False
\\ \textit{Note:} A suburb typically has more residential areas than commercial or office land use, resulting in more residents than jobs.
\item \textbf{Answer:} True
\\ \textit{Options:} A) True; B) False
\\ \textit{Note:} Centripetal forces, such as shared culture, language, and political stability, help to unify a country's citizens, fostering a sense of belonging and promoting national cohesion and stability.
\end{enumerate}

\section*{Long Form Response}
\begin{enumerate}[label=\arabic*.]
\setcounter{enumi}{10}
\item \textbf{Answer:} Primary economic activities, which involve the extraction of natural resources, are significantly influenced by the physical environment. For example, agriculture is heavily dependent on climate and soil quality; regions with fertile soil and adequate rainfall, such as the Great Plains in the United States, support extensive farming. Similarly, fishing industries thrive in coastal areas with rich marine biodiversity, like the North Atlantic Ocean, where nutrient-rich waters promote fish populations. Additionally, mining activities are often located in areas with abundant mineral deposits, such as the mining of gold in South Africa, which occurs in regions with favorable geological conditions. Overall, the physical environment shapes where and how primary economic activities develop, demonstrating the interconnectedness of geography and economy.
\\ \textit{Note:} The answer correctly identifies how primary economic activities, such as agriculture and fishing, are directly shaped by the physical environment, emphasizing the importance of climate, soil quality, and marine biodiversity in determining the viability of these industries.
\item \textbf{Answer:} The United Nations (UN) engages in interventionism during global conflicts by sending peacekeeping forces, facilitating negotiations, and providing humanitarian aid. This involvement aims to maintain peace and security, protect human rights, and help rebuild war-torn societies. By intervening, the UN can reduce violence and foster dialogue among conflicting parties, which is essential for conflict resolution. However, its actions can also lead to tensions between nations, as some countries may view intervention as an infringement on their sovereignty. Overall, the UN's interventionism plays a crucial role in shaping international relations by promoting cooperation and offering solutions to global challenges.
\\ \textit{Note:} correct because it outlines the United Nations' roles in peacekeeping, negotiation facilitation, and humanitarian assistance, emphasizing how these actions aim to reduce violence and promote dialogue, while also acknowledging the potential for international tensions arising from its interventionist policies.
\end{enumerate}

\vfill
\noindent\textit{End of Answer Key}
\end{document}